% This document provides the style to be used for a MSc Thesis at the
% Parallel and Distributed Systems group
\documentclass[11pt,twoside,a4paper,openright]{report}

% math packages
\usepackage{amsmath}
\usepackage{amssymb}

% textblocks for title page
\usepackage[absolute]{textpos}

% use babel for proper hyphenation
\usepackage[british]{babel}

% Graphics: different for pdflatex or dvi output, choose one
%%\usepackage[dvips]{graphicx}
%%\usepackage[pdftex]{graphicx}
\usepackage{graphicx}

\usepackage{epstopdf}
\usepackage{rotating}
\usepackage{subfigure}

% FONT
\usepackage[scaled=.92]{helvet}
%\usepackage{times}

% for url's use "\url{http://www.google.com/}"
\usepackage{url}
\usepackage[plainpages=false]{hyperref} 
\usepackage{cleveref}

% Information that will be filled in at various points in the report
\newcommand{\reportTitle}{Robust Downstream Communication on Transient Powered Computational RFID Devices}
\newcommand{\reportAuthor}{Jethro Tan}
\newcommand{\reportEmail}{J.E.T.Tan@student.tudelft.nl}
\newcommand{\reportUrlEmail}{\href{mailto:\reportEmail}{\reportEmail}}
\newcommand{\reportDate}{\today} %TODO: Dit is de datum van uitgifte van final versie aan de afstudeer commissie 
\newcommand{\presentationDate}{\today} %TODO: Dit is de datum van de afstudeerpresentatie 
\newcommand{\graduationCommittee}{
Prof. Dr. K. G. Langendoen & Delft University of Technology \\
Dr. Przemyslaw (Przemek) Pawelczak & Delft University of Technology \\
}

\newcommand{\reportAbstract}{TODO ABSTRACT}
\newcommand{\reportKeywords}{TODO KEYWORDS}

% For pdflatex
\pdfinfo{
   /Author (\reportAuthor)
   /Title  (\reportTitle)
   /Keywords (\reportKeywords)
}

\begin{document}

\pagenumbering{alph}
\pagestyle{empty}


% FRONTCOVER
%%\usepackage[total={210mm,297mm},left=0pt,bottom=0pt,top=0cm,right=0pt,headsep=0pt,head=0pt,showframe]{geometry}

%%\input{preambleL31958}
%%\begin{titlepage}
\begin {textblock*}{210mm}(0mm,0mm)
\noindent
\includegraphics[height=3.2cm]{pics/block}
\sffamily
\vspace{.8cm}
\begin{center}
\Large
Delft University of Technology\\
Master's Thesis in Computer Science\\
\vspace{2cm}
\parbox{170mm}{\bfseries\centering\Huge\reportTitle}\\
\vspace{1cm}
\parbox{170mm}{\bfseries\centering\reportAuthor}

\end{center}
\end{textblock*}

\begin {textblock*}{210mm}[0.0,1.0](0mm,297mm)
\noindent
\hspace{1.89cm}


\hfill\parbox{5cm}{
\includegraphics[width=5cm]{pics/es_logo_cyan_black_rgb}}
\hspace*{2cm}\\

\vspace*{1.5cm}
\noindent
\includegraphics[width=\textwidth]{pics/TU_border_A4_L_front}
\end{textblock*}

\null\newpage


%%%%%%%%%%%%%%%%%%%%%%%%%%%%%%%%%%%%%%%%%%%%%%%%%%%%%%%%%%%%%%%%%%%%%%%%%%%%%%%
\hoffset=1.63cm
\oddsidemargin=0in
\evensidemargin=0in
\textwidth=5in

%%%%%%%%%%%%%%%%%%%%%%%%%%%%%%%%%%%%%%%%%%%%%%%%%%%%%%%%%%%%%%%%%%%%%%%%%%%%%%%
\parindent=1em

% EMPTY PAGE
\cleardoublepage

\pagestyle{plain}
\pagenumbering{roman}
\setcounter{page}{1}

% TITLE PAGE: page i (hidden)
\begin{titlepage}

  \begin{center}
  \null\vfill
    \begin{center}
    \LARGE{\reportTitle}
    \end{center}

    \vspace{3cm}

    \begin{large}
    Master's Thesis in Computer Science
    \end{large}

    \vspace{1.5cm}

    \begin{normalsize}
    Embedded Software Section\\
    Faculty of Electrical Engineering, Mathematics and Computer Science\\
    Delft University of Technology\\
    Mekelweg 4, 2628 CD Delft, The Netherlands
    \end{normalsize}

    \vspace{2.0cm}

    \begin{normalsize}
    \reportAuthor \\
    \reportUrlEmail
    \end{normalsize}

    \vspace{1.0cm}

    % <MM> DD, YYYY
    \reportDate             %TODO: Dit is de datum van uitgifte van final versie aan de afstudeer commissie

  \vfill
  \end{center}

\end{titlepage}


% GRADUATION DATA AND ABSTRACT: pages ii and iii (hidden)
%De aankondiging bevat de spreker, titel, plaats, datum en tijd, samenstelling van de afstudeercommissie en een korte samenvatting (maximaal 25 regels).
\thispagestyle{empty}

\noindent \textbf{Author}\\
\begin{tabular}{l}
\reportAuthor{} (\reportUrlEmail)\\
\end{tabular}\\
\noindent \textbf{Title}\\
\begin{tabular}{l}
\reportTitle\\
\end{tabular}\\
\noindent \textbf{MSc presentation}\\
\begin{tabular}{l}
% <MM> DD, YYYY (like \today)
\presentationDate\\
\end{tabular}

\vspace{1.1cm}

\noindent \textbf{Graduation Committee}\\
\begin{tabular}{ll}
\graduationCommittee
\end{tabular}


\begin{abstract} %de abstract bevat alleen een korte samenvatting van de inhoud van het onderzoek
\setcounter{page}{3}
\reportAbstract{}
\end{abstract}

\clearpage

%\setcounter{page}{4}

% EMPTY PAGE: page iv
\cleardoublepage

% OPTIONAL QUOTATION: page v
%\include{quotation}
% EMPTY PAGE: page vi
%\cleardoublepage




% PREFACE: page v
\chapter*{Preface}
\addcontentsline{toc}{chapter}{Preface}
TODO MOTIVATION FOR RESEARCH TOPIC

\vspace{1\baselineskip}

\noindent
TODO ACKNOWLEDGEMENTS

\vspace{1\baselineskip}

\noindent
TODO AUTHOR

\vspace{1\baselineskip}

\noindent
Delft, The Netherlands

\noindent
\today


% EMPTY PAGE: page vi
\cleardoublepage


% GLOSSARY: page vii
\chapter*{Glossary}
\addcontentsline{toc}{chapter}{Glossary}








% EMPTY PAGE: page viii
\cleardoublepage



% TABLE OF CONTENTS: starting at page vii
\tableofcontents

\cleardoublepage

\pagenumbering{arabic}
\setcounter{page}{1}

% INTRODUCTION
\chapter{Introduction}
\label{intro}

\section{Computational RFID}
The WISP \cite{4539485}

\section{Challenges}

\section{Problem Statement}

\section{Contributions}
BlockWrite for the WISP
Fast downstream communication on CRFID
Protocol for wireless reprogramming of CRFID

\section{Thesis Organization}
The remainder of this thesis is as follows.
First, we introduce some the related work on 
Finally, we conclude our work in \cref{conclusions}.

% RELATED WORK
\chapter{Related Work}
\label{relatedwork}


% CONTENT CHAPTER
\chapter{Protocol Design and Implementation}
\label{design}

\section{Scribble}
The protocol was first implemented without the use of BlockWrite (an optional command in EPC C1G2).
BlockWrite was then implemented with the purpose of speeding up the downstream communication.
Both the Impinj R1000 and R420 do not execute BlockWrite as one would expect or stated in EPC C1G2 (BulkBlockWrite).

Tests were done by flashing a firmware that turned on a LED light on the WISP.
Write: ~ 4 mins
BlockWrite (using up to 5 words per BlockWrite): ~8 secs

Theoretically, one can send up to X words using BlockWrite (source?), but because Impinj performs a BulkBlockWrite rather than a BlockWrite, there is no knowing when exactly all words are received correctly.
Nor is the CRC16 that the reader sends usable verify the checksum of the whole block.
The only way is then to make use of the EPC and echo the BlockWrite buffer in it.

Alternatively, one can perform a BlockWrite with content larger than the EPC and use a Read command to verify whether all words have been written.
While using the Read to verify a BlockWrite enables us to perform larger BlockWrites, it also cuts the downstream communication time effectively in half, since the EPC is seen every round before any access command.
Furthermore, performing a larger BlockWrite also means you lose more data should the WISP come to the end of its power cycle.

% CONTENT CHAPTER
\chapter{Experimental Results}
\label{results}


% CONCLUSIONS AND FUTURE WORK
\chapter{Conclusions and Future Work}
\label{conclusions}

\section{Conclusions}

\section{Future Work}


% BIBLIOGRAPHY
\bibliographystyle{../bib/latex8}
\bibliography{../bib/article,../bib/booklet,../bib/inbook,../bib/inproceedings,../bib/misc,../bib/phdthesis,../bib/technicalreport,../bib/book,../bib/conference,../bib/incollection,../bib/manual,../bib/mscthesis,../bib/proceedings,../bib/unpublished}
\addcontentsline{toc}{chapter}{Bibliography}

% APPENDIX
%\appendix

%\include{appendix_a}

\end{document}

